\documentclass[border=0pt]{standalone}
\usepackage[RPvoltages]{circuitikz}
\usetikzlibrary{arrows,decorations.markings,arrows.meta}
\newcommand{\bushere}[3]{% length, text above, text below
    coordinate(tmp)
    % go up and do an edge down
    ++(0,#1) node[anchor=base]{#2} edge[line width=2.5pt] ++(0,{-2*#1})
    % edges do not move the current point, go down to position the node
    ++(0,{-2*#1}) node[below]{#3}
    % go back to where we started
    (tmp)
}
\begin{document}\sf
%\ctikzset{sources/fill=gray!20, resistors/fill=gray!20}
\ctikzset{sources/fill=white, resistors/fill=white}
\tikzset{myarr/.style={line width=0.74pt,-{Triangle[width=6pt, length=9pt]}}}
\begin{tikzpicture}[>=latex]% default line width
\draw [myarr](0,0) to (1.5,0);
\coordinate (B) at (1.5,0);
\coordinate (B1) at (1.5,0.05);
\path (B1)   \bushere{1}{}{};
%\draw [myarr](B)node[xshift=-4.2ex,yshift=1.3ex]{$p_i$}++(0,-0.6) 
%     to node[xshift=-0.5ex,yshift=1.3ex]{$p_o$} ++(1.5,0);

%\draw [myarr](B)node[xshift=-4.2ex,yshift=1.3ex]{$p_i$}++(0,0) 
%     to node[xshift=-0.5ex,yshift=1.3ex]{$p_o$} ++(1.5,0);
\draw [myarr](B)node[xshift=-5ex,yshift=1.3ex]{$p_i$}++(0,0.6) 
     to node[xshift=-0.5ex,yshift=1.3ex]{$p_d$} ++(1.5,0);
%\draw [myarr](B)node[xshift=-4.2ex,yshift=1.3ex]{$p_i$}++(0,0.9) 
%     to node[xshift=-0.5ex,yshift=1.3ex]{$p_o$} ++(1.5,0);
%\draw [myarr](B)node[xshift=-4.2ex,yshift=1.3ex]{$p_i$}++(0,-0.9) 
%     to node[xshift=-0.5ex,yshift=1.3ex]{$p_o$} ++(1.5,0);




%\draw (B)to ++(1,1);
%\draw (B)to ++(1,-1);

\end{tikzpicture}
\end{document}